\chapter{Omtrek en oppervlakte}\label{ch:g5:areas}

In jaar 4 telden we rastervakjes (\cref{ch:g4:measure}); in dit hoofdstuk
verdienen we de eerste \emph{formule} voor de \omterm{def:g4:measure:area}{oppervlakte} --- lengte keer
breedte voor de rechthoek --- en leren we de eenheden cm$^2$ en m$^2$. De
formules voor driehoeken en andere figuren rijpen in \cref{ch:g6:measure} en
\cref{ch:g7:areas}.

\section{Twee verschillende maten}

\begin{example}[Rand tegenover oppervlak]\label{ex:g5:areas:contrast}
Een tuinier heeft de \emph{\omterm{def:g4:measure:perimeter}{omtrek}} nodig om hek te kopen en de
\emph{\omterm{def:g4:measure:area}{oppervlakte}} om graszaad te kopen. De twee volgen elkaar niet: een
rechthoek dunner en langer maken kan zijn \omterm{def:g4:measure:area}{oppervlakte} gelijk houden terwijl de
\omterm{def:g4:measure:perimeter}{omtrek} groeit --- vergelijk een rechthoek van $6 \times 2$ met een van
$4 \times 3$ (dezelfde \omterm{def:g4:measure:area}{oppervlakte} $12$, omtrekken $16$ en $14$).
\end{example}

\section{De formule voor de rechthoek}

\begin{proposition}[Oppervlakte van een rechthoek]\label{prop:g5:areas:rectangle}
Een rechthoek met lengte $l$ en breedte $b$ (in dezelfde eenheid) heeft
\omterm{def:g4:measure:area}{oppervlakte}
\[
A = l \times b ,
\]
in vierkante eenheden: cm$^2$ (vierkantjes met zijde $1$ cm), m$^2$
(vierkanten met zijde $1$ m), \dots
\end{proposition}

\begin{proof}[Waarom, door te tellen]
Bedek de rechthoek met eenheidsvierkantjes: $b$ rijen van $l$ vierkantjes,
dus $l \times b$ vierkantjes in totaal --- precies de
vermenigvuldigrechthoek van \cref{ch:g3:mult}.
\end{proof}

\begin{omfigure}
\begin{tikzpicture}[scale=0.62]
  \draw[gray!50, thin] (0,0) grid (6,3);
  \draw[very thick, omDef] (0,0) rectangle (6,3);
  \fill[omThm!30] (0,2) rectangle (1,3);
  \draw[omThm, thick] (0,2) rectangle (1,3);
  \node[right, font=\small, omThm] at (1.15,2.5) {$1$ cm$^2$};
  \node[below, font=\small] at (3,-0.15) {$6$ cm};
  \node[left, font=\small] at (0,1.5) {$3$ cm};
  \node[font=\small] at (3.9,1.2) {$6 \times 3 = 18$ cm$^2$};
\end{tikzpicture}

{\small Drie rijen van zes centimetervierkantjes: de formule $l \times b$ is
dat tellen, eens en voor altijd opgeschreven.}
\end{omfigure}

\begin{example}\label{ex:g5:areas:rectangle}
Een vloerkleed meet $2.5$ m bij $2$ m: \omterm{def:g4:measure:area}{oppervlakte} $2.5 \times 2 = 5$ m$^2$.
Een postzegel meet $3$ cm bij $2.4$ cm: \omterm{def:g4:measure:area}{oppervlakte} $3 \times 2.4 = 7.2$
cm$^2$. Dezelfde formule, elke eenheid --- zolang beide zijden \emph{dezelfde}
eenheid gebruiken.
\end{example}

\begin{example}[Vierkante eenheden gaan per 100]\label{ex:g5:areas:units}
$1$ m $= 100$ cm, maar $1$ m$^2 = 100 \times 100 = 10\,000$ cm$^2$: een
vierkante meter is een raster van $100$ bij $100$ vierkante centimeters.
Oppervlakte-eenheden springen per honderd, niet per tien.
\end{example}

\section{Samengestelde figuren}

\begin{method}[Knippen, optellen, aftrekken]\label{met:g5:areas:composite}
Voor een figuur die uit rechthoeken bestaat (een L, een T, een lijst):
\begin{enumerate}
  \item knip haar in rechthoeken, of vul haar aan tot één grote rechthoek;
  \item reken elke rechthoekige \omterm{def:g4:measure:area}{oppervlakte} met de formule uit;
  \item tel de stukken op --- of trek het gat van de grote rechthoek af;
  \item controleer met een ruwe telling van rastervakjes.
\end{enumerate}
\end{method}

\begin{example}\label{ex:g5:areas:composite}
Een L-vormige kamer: een rechthoek van $6$ m $\times$ $4$ m met een hoek van
$2$ m $\times$ $2$ m eruit.
\[
\text{Oppervlakte} = 6 \times 4 - 2 \times 2 = 24 - 4 = 20 \text{ m}^2 .
\]
Of knip de L in een strook $6 \times 2$ en een strook $4 \times 2$:
$12 + 8 = 20$ m$^2$ --- twee wegen, één antwoord.
\end{example}

\begin{example}[De helft van een rechthoek]\label{ex:g5:areas:halfrect}
Een rechthoek langs een diagonaal knippen geeft twee driehoeken met dezelfde
\omterm{def:g4:measure:area}{oppervlakte}: elk is de \emph{\omterm{def:g3:division:half}{helft}} van de rechthoek. Een rechthoekige
driehoek met rechthoekszijden $6$ cm en $4$ cm heeft dus \omterm{def:g4:measure:area}{oppervlakte}
$\frac{6 \times 4}{2} = 12$ cm$^2$ --- een plaatje dat het waard is om te
onthouden voor \cref{ch:g6:measure}, waar het een formule wordt.
\end{example}

\begin{omfigure}
\begin{tikzpicture}[scale=0.62]
  \draw[gray!50, thin] (0,0) grid (6,4);
  \draw[very thick] (0,0) rectangle (6,4);
  \fill[omDef!25] (0,0) -- (6,0) -- (0,4) -- cycle;
  \draw[very thick, omDef] (6,0) -- (0,4);
  \node[font=\small] at (1.7,1.2) {helft};
  \node[font=\small] at (4.3,2.8) {helft};
\end{tikzpicture}

{\small De diagonaal knipt de rechthoek van $6 \times 4$ in twee gelijke
driehoeken van $12$ vakjes elk.}
\end{omfigure}

\section{Oefeningen}

\begin{exercise}[$\star$]\label{exo:g5:areas:1}
Reken de \omterm{def:g4:measure:area}{oppervlakte} en de \omterm{def:g4:measure:perimeter}{omtrek} uit van een rechthoek van $8$ cm $\times$
$5$ cm. Welk antwoord is in cm en welk in cm$^2$?
\end{exercise}

\begin{exercise}[$\star$]\label{exo:g5:areas:2}
Reken de \omterm{def:g4:measure:area}{oppervlaktes} uit: een vierkant met zijde $9$ cm; een rechthoek van
$12$ m $\times$ $7$ m; een rechthoek van $4.5$ cm $\times$ $6$ cm.
\end{exercise}

\begin{exercise}[$\star$]\label{exo:g5:areas:3}
Een rechthoek heeft \omterm{def:g4:measure:area}{oppervlakte} $63$ cm$^2$ en lengte $9$ cm. Zoek zijn
breedte, en daarna zijn \omterm{def:g4:measure:perimeter}{omtrek}.
\end{exercise}

\begin{exercise}[$\star$]\label{exo:g5:areas:4}
Teken op ruitjespapier twee verschillende rechthoeken met een \omterm{def:g4:measure:area}{oppervlakte} van
$24$ vakjes, en reken beide omtrekken uit. Dezelfde \omterm{def:g4:measure:area}{oppervlakte} --- dezelfde
\omterm{def:g4:measure:perimeter}{omtrek}?
\end{exercise}

\begin{exercise}[$\star$]\label{exo:g5:areas:5}
Zet om: $3$ m$^2$ in cm$^2$; \quad $50\,000$ cm$^2$ in m$^2$. (Denk aan
\cref{ex:g5:areas:units}: per honderd!)
\end{exercise}

\begin{exercise}[$\star$]\label{exo:g5:areas:6}
Een T-vormige figuur bestaat uit een horizontale balk van $8 \times 2$ boven
een verticale balk van $2 \times 5$ (in cm). Reken haar \omterm{def:g4:measure:area}{oppervlakte} uit door
twee rechthoeken op te tellen.
\end{exercise}

\begin{exercise}[$\star$]\label{exo:g5:areas:7}
Een fotolijst: een rechthoek van $30$ cm $\times$ $20$ cm met een
rechthoekige opening van $24$ cm $\times$ $14$ cm eruit. Welke \omterm{def:g4:measure:area}{oppervlakte} aan
hout gebruikt de lijst?
\end{exercise}

\begin{exercise}[$\star$]\label{exo:g5:areas:8}
Reken de \omterm{def:g4:measure:area}{oppervlakte} uit van een rechthoekige driehoek met rechthoekszijden
$8$ cm en $6$ cm (de \omterm{def:g3:division:half}{helft} van een rechthoek,
\cref{ex:g5:areas:halfrect}).
\end{exercise}

\begin{exercise}[$\star$]\label{exo:g5:areas:9}
Een rechthoekig tuinbed meet $7$ m bij $4$ m. Zaad kost $2$ per vierkante
meter en hek $3$ per meter. Reken de kosten van het zaad uit, en daarna die
van het hek.
\end{exercise}

\begin{exercise}[$\star\star$]\label{exo:g5:areas:10}
De vloer van een gang is $12$ m lang en $2$ m breed en moet met vierkante
tegels met zijde $50$ cm bedekt worden. Hoeveel tegels zijn er nodig? (Zet
eerst om, of tel de tegels in elke richting.)
\end{exercise}

\begin{exercise}[$\star\star$]\label{exo:g5:areas:11}
Verdubbel de zijden van een rechthoek van $4$ cm $\times$ $3$ cm. Wat gebeurt
er met zijn \omterm{def:g4:measure:perimeter}{omtrek}? En met zijn \omterm{def:g4:measure:area}{oppervlakte}? (Reken beide voor en na uit ---
de twee antwoorden verschillen!)
\end{exercise}
